\section{Data \& Synthetic Sampling}
\label{sec:data}

\subsection{Brain Surface Mesh and Graph Construction}

We use midthickness cortical surface meshes from the Human Connectome Project (HCP)~\cite{vanessen2013hcp},
consisting of approximately 32{,}492 vertices per hemisphere.
For each subject we extract the left hemisphere and restrict our analysis to an occipital patch
defined as all vertices within a geodesic radius of 40\,mm from the occipital pole, a threshold
chosen to approximately capture the extent of primary visual cortex (V1)~\cite{wandell2007}.

The occipital pole is identified as the most posterior vertex (minimum $y$-coordinate in MNI space)
and serves as the foveal anchor point, consistent with the observation that the foveal representation
lies just anterior to the occipital pole~\cite{sereno1995borders}.

From the triangulated mesh we build an undirected graph $G = (V, E)$ where vertices $V$ correspond to
cortical nodes and edges $E$ connect nodes sharing a triangle edge.
Edge weights are set to the Euclidean distance between the 3D coordinates of connected nodes (in mm).
Geodesic surface distances are computed via Dijkstra's algorithm on this weighted graph.

\subsection{Synthetic Retinotopic Map}

\paragraph{Eccentricity.}
The geodesic distance $d_\text{geo}$ from each node to the occipital pole is converted to
visual field eccentricity using a cortical magnification law of the form $D(r) = A\log(r+B)$~\cite{schwartz1980}.
We set $B = 1.6$, following the human magnification model of~\citet{drasdo1977} as cited by~\citet{schwartz1980},
and $A = 15.1\,\text{mm/deg}$, the foveal magnification factor reported by~\citet{cowey1974}:
\begin{equation}
  \varepsilon = \exp\!\left(\frac{d_\text{geo}}{15.1}\right) - 1.6,
  \quad \varepsilon \in [0.1^\circ,\, 12^\circ].
\end{equation}

\paragraph{Polar angle.}
The polar angle $\theta$ encodes the compass direction of each node around the occipital pole,
following the concept of polar-angle retinotopic mapping established by phase-encoded fMRI~\cite{sereno1995borders},
computed here in the $x$-$z$ plane of MNI space:
\begin{equation}
  \theta_i = \arctan2\!\left(z_i - z_\text{pole},\; x_i - x_\text{pole}\right).
\end{equation}
The $x$-$z$ plane of the brain approximates the visual field plane for the left hemisphere,
so $\theta$ corresponds to the direction in the right visual hemifield.

\paragraph{Visual field coordinates.}
Eccentricity and polar angle are converted to Cartesian visual field coordinates:
\begin{equation}
  \mathrm{vf\_x}_i = \varepsilon_i \cos\theta_i, \qquad
  \mathrm{vf\_y}_i = \varepsilon_i \sin\theta_i.
\end{equation}

\subsection{pRF Signal Simulation}

We simulate population receptive field (pRF) responses following the Gaussian pRF model~\cite{dumoulin2008}.
For each of $S = 100$ randomly placed stimuli, the response of node $i$ is:
\begin{equation}
  r_{i,s} = \exp\!\left(-\frac{1}{2}\left(\frac{d_{i,s}}{\sigma_i}\right)^{\!2}\right),
\end{equation}
where $d_{i,s} = \|\mathrm{vf}_i - \mathrm{stim}_s\|_2$ is the distance between the node's preferred
visual field position and stimulus $s$ (in degrees), and $\sigma_i$ is the pRF size:
\begin{equation}
  \sigma_i = 0.5 + 0.15\,\varepsilon_i.
\end{equation}
The linear scaling of $\sigma$ with eccentricity reflects cortical magnification:
peripheral nodes have larger receptive fields and thus respond to a broader range of stimulus positions,
consistent with the V1 pRF size-eccentricity relationship reported by~\citet{harvey2011}
(intercept and slope derived from their reported pRF size at $4^\circ$ eccentricity and slope-to-intercept ratio).
The resulting feature matrix $X \in \mathbb{R}^{N \times 100}$ forms the GNN input.

